\subsection{Aims and objectives}
\label{sec:aims}

This project develops a protocol-aware Deep Reinforcement Learning (DRL) framework that
optimises UAV flight trajectories for energy-efficient data collection in dense, unsynchronised
LoRaWAN networks. The central claim is that using the UAV's physical position as a direct
control input for the radio protocol resolves the ``Spreading Factor (SF) Monoculture''
bottleneck. A Deep Q-Network (DQN) agent is trained within a high-fidelity simulation to
learn a network-aware policy that maximises data throughput while minimising energy consumption.

The project objectives are:

\begin{itemize}
    \item To develop a high-fidelity simulation environment compatible with the Gymnasium
    interface, incorporating rigorous telecommunications physics including the Two-Ray Ground
    Reflection path loss model, log-normal shadowing, and asynchronous duty cycles. To the
    author's knowledge, this is the first open framework to model the complete causal chain
    from UAV position through RSSI, EMA-ADR, and Capture Effect to packet delivery within a
    single simulation step.

    \item To model the specific constraints of the LoRaWAN protocol, particularly the
    Adaptive Data Rate (ADR) latency, to replicate the causal link between UAV positioning,
    RSSI, and Spreading Factor allocation.

    \item To formulate the UAV data collection problem as a Partially Observable Markov
    Decision Process (POMDP, where sensor locations and out-of-range channel states are
    hidden from the agent), and to approximate the resulting POMDP as a tractable MDP via
    frame stacking ($k{=}4$), designing a multi-objective reward function that balances the
    competing goals of energy conservation, data freshness (AoI), and collection volume.

    \item To train a Deep Q-Network (DQN) agent to learn an emergent flight policy that
    autonomously navigates the ``Starvation vs.\ Throughput'' trade-off without human
    intervention.

    \item To evaluate the protocol-aware DRL agent against two greedy path-planning baselines
    (Nearest-Sensor Greedy and SF-Aware Max-Throughput Greedy) in terms of energy efficiency,
    fairness (Jain's Index), and total bytes delivered, using multi-seed statistical testing
    (Welch's $t$-test, Cohen's $d$).

    \item To characterise the operational envelope of the single-UAV system by evaluating
    across grid scales from $100{\times}100$ to $1{,}000{\times}1{,}000$ cells, identifying
    the physical feasibility boundary beyond which link-budget and energy constraints prevent
    full sensor coverage regardless of policy quality.

    \item To assess the trade-offs between computational complexity, policy robustness, and
    real-world feasibility, identifying the ``Sim-to-Real'' gap and the transition from
    feasible to infeasible regimes as key findings for future research.
\end{itemize}

A protocol-aware agent generates energy-efficient trajectories satisfying network fairness and throughput requirements without pre-specified routing heuristics.
